% Stage 8 restored fragment — 2.2.3 Causality and time-frequency equivalence (archive lines 3006-3480)
\subsection{Causality and time-frequency equivalence}
\label{sec:ch223}

A standard-gain horn on a compact-range turntable fixes the laboratory scene for `Causality and time-frequency equivalence'': frequency, probe standoff, and wattmeter closure on $\Gamma_R$ precede the formal displays below \cite{hansen1988}. Maxwell's retarded export to $\Gamma_R$ remains the governing constraint \cite{maxwell1873}.
The physical problem is retarded response: no field at $(\mathbf{r},t)$ may depend on source
events that have not yet had time to propagate from $\mathbf{r}'$ to $\mathbf{r}$
\cite{maxwell1873}. On the range, this is the difference between a
measured impulse response that vanishes before the light-travel delay and an unphysical
advanced template that leads the feed. A pulsed compact-range measurement makes the
distinction visible: energy arrives only after transit time $|\mathbf{r}-\mathbf{r}'|/c_0$;
a single-frequency phasor can hide violations that broadband synthesis would expose
\cite{harrington2001}. Causality is therefore the principle that links time-domain evolution
to harmonic analysis before any Laplace or resolvent notation is introduced
\cite{stratton1941}. A time-domain operator is causal when no response can occur before the
corresponding excitation \cite{stratton1941}. In electromagnetics this requirement appears
through retarded propagation: disturbances generated at one location influence another only
after the finite transit time imposed by the speed of light \cite{sommerfeld1949}.
Subsection~\ref{sec:ch223} supplies the definitive proof that $\mathcal{S}_\omega$ and
$\mathcal{R}_\omega$ are the frequency-domain images of causal time-domain maps when
impressed data, media, and boundary class are held fixed \cite{stratton1941}. The retarded
Eq.~\eqref{eq:ch2.causality-retarded} and symbolic relation
Eq.~\eqref{eq:ch2.time-to-frequency-map} are upgraded to an equivalence theorem;
Green-function construction remains in Section~\ref{sec:ch23} \cite{stratton1941}.
Causal support. A time-domain operator \(\mathcal{K}\) is \emph{causal} if

\begin{equation}
\label{eq:ch2.causal-operator-support}
\mathcal{K}[f](t)=0\ \text{for}\ t<t_0
\quad\text{whenever}\ 
\operatorname{supp}_t f\subseteq[t_0,\infty),
\end{equation}
for admissible test functions \(f\) \cite{stratton1941}. Retarded propagation forbids
response before transit time $|\mathbf{r}-\mathbf{r}'|/c$; frequency-domain Green
functions used on ranges are analytic continuations of that support, while broadband
pulses expose violations that single-frequency phasors can hide \cite{sommerfeld1949}.
The retarded dyadic kernel
in Eq.~\eqref{eq:ch2.retarded-convolution-form} satisfies
\begin{equation}
\label{eq:ch2.retarded-kernel-support}
\mathcal{G}(\mathbf{r},\mathbf{r}',t-t')=\mathbf{0}\quad\text{for}\ t-t'<0,
\end{equation}
which enforces Eq.~\eqref{eq:ch2.causal-support-constraint} on \(\E\)
\cite{stratton1941}. Advanced kernels with
\(\mathcal{G}(\mathbf{r},\mathbf{r}',t-t')\neq\mathbf{0}\) for \(t-t'<0\) violate
Eq.~\eqref{eq:ch2.causal-operator-support} and are excluded from the equivalence
class treated here.
\begin{figure}[t]
\centering
\ChOneFigBlock{%
\begin{tikzpicture}[ch1fig, node distance=5mm and 7mm]
 \node[ch1box] (td) {$\mathcal{S}_t$\\time domain};
 \node[ch1box, right=7mm of td] (lt) {Laplace\\$s=-\jj\omega$};
 \node[ch1box, right=7mm of lt] (fd) {$\mathcal{S}_\omega$\\harmonic};
 \draw[ch1flow] (td) -- (lt) -- (fd);
\end{tikzpicture}%
}
\caption{Causality link: causal \(\mathcal{S}_t\) maps to harmonic
\(\mathcal{S}_\omega\) under Laplace/Fourier evaluation on the locked phasor axis
\(s=-\jj\omega\), \(\omega>0\).}
\label{fig:ch2.laplace-resolvent-linkage}
\end{figure}
Let
\(\mathbf{j}_{\mathrm{imp}}(\mathbf{r},t)\) be an impressed current with compact
support in \(\Omega_s\times[0,\infty)\) and let \(\mathbf{U}_0=\mathbf{0}\) for
simplicity. The Laplace transform in \(s\) with \(\Re\{s\}\) sufficiently large is
\begin{equation}
\label{eq:ch2.laplace-source-transform}
\widetilde{\mathbf{j}}(\mathbf{r},s)
=
\int_0^\infty \mathbf{j}_{\mathrm{imp}}(\mathbf{r},t)\,e^{-st}\,dt.
\end{equation}
On the phasor axis \(s=-\jj\omega\) with \(\omega>0\),
\begin{equation}
\label{eq:ch2.fourier-phasor-identification}
\widetilde{\mathbf{j}}(\mathbf{r},-\jj\omega)
=
\int_0^\infty \mathbf{j}_{\mathrm{imp}}(\mathbf{r},t)\,e^{\jj\omega t}\,dt,
\end{equation}
which is the Fourier transform restricted to causal support
\cite{stratton1941}. For \(\mathbf{j}_{\mathrm{imp}}(\mathbf{r},t)=\Re\{\underline{\mathbf{j}}(\mathbf{r})e^{-\jj\omega_0 t}\}\)
with fixed \(\omega_0>0\), Eq.~\eqref{eq:ch2.fourier-phasor-identification}
selects the phasor \(\underline{\mathbf{j}}\) under the locked convention
\(\exp(-\jj\omega t)\) \cite{stratton1941}. Applying the Laplace transform to Eq.~\eqref{eq:ch2.time-ibvp-package} with fixed
\(\mathcal{B}_\Gamma\) and homogeneous initial data yields, formally,
\begin{equation}
\label{eq:ch2.laplace-maxwell-resolvent}
\widetilde{\mathbf{U}}(\mathbf{r},s)
=
\mathcal{R}(s)\,\widetilde{\mathbf{F}}(s),
\qquad
\mathcal{R}(s)=\mathcal{L}_s^{-1},
\end{equation}
where \(\mathcal{L}_s\) is the Laplace-domain Maxwell operator obtained from
\(\mathcal{A}\) by Eq.~\eqref{eq:ch2.frequency-substitution-rule} with
\(\partial_t\to s\) \cite{stratton1941}. On the axis \(s=-\jj\omega\),
\(\mathcal{L}_s=\mathcal{L}_\omega\) from Eq.~\eqref{eq:ch2.composite-operator-domain},
so
\begin{equation}
\label{eq:ch2.resolvent-green-equivalence}
\mathcal{R}(-\jj\omega)=\mathcal{L}_\omega^{-1}
\quad\text{on the outgoing/resolvent branch,}
\end{equation}
and the retarded convolution Eq.~\eqref{eq:ch2.retarded-convolution-form} is the
inverse Laplace transform of Eq.~\eqref{eq:ch2.laplace-maxwell-resolvent} when
\(\mathcal{G}\) is the causal Green operator constructed in Section~\ref{sec:ch23}
\cite{stratton1941}. Equation~\eqref{eq:ch2.resolvent-green-equivalence}
is the operator form of Eq.~\eqref{eq:ch2.causality-retarded}: harmonic
inversion is not an independent postulate but the \(s=-\jj\omega\) evaluation of a
causal resolvent.
The equivalence theorem requires explicit assumptions:
\begin{equation}
\label{eq:ch2.equivalence-assumptions}
\begin{aligned}
&\text{(i)}\ \mathcal{B}_\Gamma\ \text{and}\ \varepsilon,\mu\ \text{fixed in time},\\
&\text{(ii)}\ \mathbf{j}_{\mathrm{imp}}\ \text{causal with}\ \operatorname{supp}_t\subseteq[0,\infty),\\
&\text{(iii)}\ \text{outgoing/resolvent branch selected for exterior domains},\\
&\text{(iv)}\ \mathbf{U}_0=\mathbf{0}\ \text{or compatible with harmonic steady limit}.
\end{aligned}
\end{equation}
Under Eqs.~\eqref{eq:ch2.equivalence-assumptions},
\begin{equation}
\label{eq:ch2.harmonic-equivalence-theorem}
\mathcal{S}_\omega[\underline{\mathbf{j}}]
=
\eval{\widetilde{\mathbf{U}}(\mathbf{r},s)}_{s=-\jj\omega}
=
\eval{\mathcal{L}_\omega^{-1}\widetilde{\mathbf{F}}(s)}_{s=-\jj\omega},
\end{equation}
where \(\underline{\mathbf{j}}\) is the phasor of
\(\mathbf{j}_{\mathrm{imp}}(\mathbf{r},t)=\Re\{\underline{\mathbf{j}}(\mathbf{r})e^{-\jj\omega t}\}\)
and \(\widetilde{\mathbf{F}}\) is the Laplace forcing induced by
\(\widetilde{\mathbf{j}}\) \cite{stratton1941}. Equation
\eqref{eq:ch2.harmonic-equivalence-theorem} is the formal statement that
frequency-domain solvers implement causal time-domain maps when
Eqs.~\eqref{eq:ch2.equivalence-assumptions} hold \cite{harrington2001}. For
single-frequency excitation with damping or outgoing radiation,
\begin{equation}
\label{eq:ch2.steady-state-harmonic-limit}
\lim_{t\to\infty}\mathbf{U}(t)
=
\Re\big\{\mathcal{S}_\omega[\underline{\mathbf{j}}]\,e^{-\jj\omega t}\big\},
\end{equation}
which justifies steady-state phasor measurements after switch-on transients decay
\cite{stratton1941}. Because $\Pi_{\mathrm{rad}}$ is linear and commutes with harmonic
extraction on the outgoing subspace,
\begin{equation}
\label{eq:ch2.radiation-time-frequency-equiv}
\mathcal{R}_\omega[\underline{\mathbf{j}}]
=
\Pi_{\mathrm{rad}}\,\mathcal{S}_\omega[\underline{\mathbf{j}}],
\end{equation}
with transport and regime qualifications from
Eq.~\eqref{eq:ch2.transport-qualified-radiation} unchanged
\cite{stratton1941}. Frequency-domain codes that invert $\mathcal{L}_\omega$ without
outgoing or passive completion implement a different resolvent branch than the causal
retarded $\mathcal{G}$ of Eq.~\eqref{eq:ch2.retarded-convolution-form}
\cite{harrington2001}. Advanced Green functions satisfy
\begin{equation}
\label{eq:ch2.advanced-kernel-violation}
\mathcal{G}_{\mathrm{adv}}(\mathbf{r},\mathbf{r}',t-t')\neq\mathbf{0}
\ \text{for some}\ t-t'<0,
\end{equation}
and cannot satisfy Eq.~\eqref{eq:ch2.harmonic-equivalence-theorem} for any causal
$\mathbf{j}_{\mathrm{imp}}$ \cite{stratton1941}. Broadband synthesis requires
superposition over $\omega$, and time-varying $\mathcal{B}_\Gamma(t)$ breaks assumption
(i) in Eq.~\eqref{eq:ch2.equivalence-assumptions} \cite{stratton1941}.
Subsections~\ref{sec:ch221}--\ref{sec:ch223} fixed evolution, harmonic operators, and
causality. Subsection~\ref{sec:ch224} supplies spectral partitions into propagating,
evanescent, and non-propagating subspaces in operator form \cite{stratton1941}.
Reactive versus radiative regime criteria from Section~\ref{sec:ch01.2.4} are invoked by
reference; the outgoing selector $\Pi_{\mathrm{rad}}$ of
Eq.~\eqref{eq:ch2.radiation-operator-def} is refined spectrally, not replaced
\cite{poynting1884}.
The physical problem is spectral transport: a near-field scanner records both traveling and
decaying longitudinal content, yet only outgoing propagating modes may enter far-zone power
reports \cite{hansen1988}. Subsection~\ref{sec:ch224} defines propagating, evanescent, and non-propagating spectral
subspaces and refines \(\Pi_{\mathrm{rad}}=\mathcal{O}_{\mathrm{out}}\circ\Pi_{\mathrm{prop}}\)
so that only outgoing propagating content qualifies for far-zone transport reports
\cite{stratton1941}. Recall that $\mathcal{S}_\omega:\mathcal{J}_\Omega\to\mathcal{F}_\Omega$ produces the
full harmonic field pair, $\mathcal{R}_\omega=\Pi_{\mathrm{rad}}\circ\mathcal{S}_\omega$
retains only radiative content, and $\Pi_{\mathrm{rad}}$ acts after propagating selection
\cite{harrington2001}. A near-field scanner routinely measures both propagating and evanescent spectral content.
The distinction becomes visible when reconstructed fields are propagated away from the
measurement surface: propagating components persist, whereas evanescent components decay
exponentially with distance \cite{hansen1988}. The spectral decomposition
developed below states that experimentally familiar observation in spectral language.
In a homogeneous region with symbol \(k_0\) from Eq.~\eqref{eq:ch2.helmholtz-symbol},
plane-wave trials decompose the field into lateral and longitudinal wavenumber components.
The dispersion relation below is the local criterion that
separates traveling from decaying longitudinal behavior: when lateral wavenumber
\(k_t=\sqrt{k_x^2+k_y^2}\) exceeds \(k_0\), no real propagating longitudinal root exists
and the mode becomes evanescent \cite{stratton1941}. Experimentally one would observe that
near-field scan content with \(k_t>k_0\) dies away within a penetration depth
\(\delta_{\mathrm{ev}}=1/\alpha(k_t)\) as the field is propagated outward
\cite{hansen1988}. Plane-wave trials \(\mathbf{u}_{\mathbf{k}}(\mathbf{r})=\mathbf{u}_0\,e^{-\jj\mathbf{k}\cdot\mathbf{r}}\)
satisfy the Helmholtz condition when
\begin{equation}
\label{eq:ch2.dispersion-relation}
|\mathbf{k}|=k_0,\qquad
k_z^2=k_0^2-k_x^2-k_y^2,
\end{equation}
with lateral wavenumber \(k_t=\sqrt{k_x^2+k_y^2}\) \cite{stratton1941}. Equation
\eqref{eq:ch2.dispersion-relation} splits spectral content by longitudinal behavior:
\begin{equation}
\label{eq:ch2.longitudinal-wavenumber-split}
k_z=
\begin{cases}
\pm\sqrt{k_0^2-k_t^2}, & k_t\le k_0\ \text{(propagating)},\\[4pt]
-\jj\,\alpha(k_t), & k_t>k_0\ \text{(evanescent)},\ \alpha>0,
\end{cases}
\end{equation}
where the evanescent root is chosen with decay in the \(+z\) half-space
\cite{stratton1941}. For \(k_t>k_0\) the decay rate is
\begin{equation}
\label{eq:ch2.evanescent-decay-rate}
\alpha(k_t)=\sqrt{k_t^2-k_0^2},
\end{equation}
so a pure evanescent trial on \(z>0\) scales as
\(e^{-\alpha(k_t)z}\) with penetration depth \(\delta_{\mathrm{ev}}=1/\alpha(k_t)\)
\cite{stratton1941}. Propagating modes carry real \(k_z=\pm\sqrt{k_0^2-k_t^2}\) and
oscillatory phase along the longitudinal axis; evanescent modes carry imaginary
\(k_z\) and store reactive energy in the near zone without far-zone power export
\cite{harrington2001}. At fixed \((k_x,k_y)\), the Helmholtz symbol
\(\mathcal{L}_\omega\) from Subsection~\ref{sec:ch222} acts on each longitudinal
branch as a scalar multiplier in the plane-wave basis, which is why spectral
partitions commute with harmonic extraction once outgoing radiation condition is fixed on the
propagating sheet \cite{stratton1941}.
Transverse polarization completes the mode description. For a trial electric field
\(\widetilde{\E}(k_x,k_y)\,e^{-\jj(k_x x+k_y y+k_z z)}\), source-free admissibility
requires \(\mathbf{k}\cdot\widetilde{\E}=0\) in the homogeneous region
\cite{stratton1941}. Decompose \(\mathbf{k}=\mathbf{k}_t+k_z\hat{\mathbf{z}}\)
with \(\mathbf{k}_t=k_x\hat{\mathbf{x}}+k_y\hat{\mathbf{y}}\). Then
\begin{equation}
\label{eq:ch2.transverse-polarization-constraint}
k_x\widetilde{E}_x+k_y\widetilde{E}_y+k_z\widetilde{E}_z=0,
\end{equation}
which leaves two independent transverse amplitudes on each branch of
Eq.~\eqref{eq:ch2.longitudinal-wavenumber-split}. On the propagating sheet,
Eq.~\eqref{eq:ch2.transverse-polarization-constraint} supports traveling power
along \(\hat{\mathbf{z}}\) when outgoing \(k_z\) is selected; on the evanescent sheet
the same constraint forces predominantly tangential storage with longitudinal
\(E_z\) tied to \(\alpha(k_t)\) \cite{stratton1941}. Subsection~\ref{sec:ch224}
does not classify TE/TM guide nomenclature---that labeling is deferred to
waveguide sections---but Eq.~\eqref{eq:ch2.transverse-polarization-constraint}
is the local polarization criterion paired with the dispersion criterion of
Eq.~\eqref{eq:ch2.dispersion-relation} \cite{harrington2001}.
\begin{figure}[t]
\centering
\ChOneFigBlock{%
\begin{tikzpicture}[ch1fig, node distance=5mm and 7mm]
 \node[ch1box] (kt) {$k_t$};
 \node[ch1box, right=7mm of kt] (prop) {$k_t\le k_0$\\propagating};
 \node[ch1box, below=5mm of kt] (ev) {$k_t>k_0$\\evanescent};
 \draw[ch1flow] (kt.east) -- (kt.east |- prop.west);
 \draw[ch1flow] (kt.east |- prop.west) -- (prop.west);
 \draw[ch1flow] (kt.south) -- (kt.south |- ev.north);
 \draw[ch1flow] (kt.south |- ev.north) -- (ev.north);
\end{tikzpicture}%
}
\caption{Dispersion split Eq.~\eqref{eq:ch2.longitudinal-wavenumber-split}: lateral
wavenumber \(k_t\) selects propagating versus evanescent longitudinal behavior at
fixed \(k_0\).}
\label{fig:ch2.propagating-evanescent-split}
\end{figure}
Angular-spectrum representations in
Section~\ref{sec:ch29} expand fields in these modes; here the subspaces are defined
as solution domains \cite{stratton1941}. Let \(\mathcal{F}_{\mathrm{prop}}\subset\mathcal{F}_\Omega\) be the completion of
finite superpositions of propagating plane waves with outgoing \(k_z\) selected by
exterior completion, let \(\mathcal{F}_{\mathrm{ev}}\subset\mathcal{F}_\Omega\) be the
closure of evanescent superpositions, and let
\(\mathcal{F}_{\mathrm{non}}\subset\mathcal{F}_\Omega\) collect modes that are
neither outgoing-propagating nor evanescent-decaying in the declared half-space---for
example cutoff guide modes and purely reactive standing content tied to terminals
\cite{stratton1941}. 
\phantomsection\label{sec:ch224}
Orthogonal spectral projectors (formal) satisfy

\begin{equation}
\label{eq:ch2.spectral-projectors}
\Pi_{\mathrm{prop}}+\Pi_{\mathrm{ev}}+\Pi_{\mathrm{non}}=\mathrm{id}
\quad\text{on}\ \widehat{\mathcal{F}}_\Omega,
\end{equation}
with \(\widehat{\mathcal{F}}_\Omega\) the spectral completion of admissible states
\cite{stratton1941}. Idempotency and mutual annihilation read
\begin{equation}
\label{eq:ch2.spectral-projector-algebra}
\Pi_{\mathrm{prop}}^2=\Pi_{\mathrm{prop}},\quad
\Pi_{\mathrm{ev}}^2=\Pi_{\mathrm{ev}},\quad
\Pi_{\mathrm{non}}^2=\Pi_{\mathrm{non}},\quad
\Pi_{\mathrm{prop}}\Pi_{\mathrm{ev}}=\Pi_{\mathrm{ev}}\Pi_{\mathrm{prop}}=0,
\end{equation}
with analogous zero products for \(\Pi_{\mathrm{non}}\) \cite{stratton1941}. On
\(\mathcal{D}(\mathcal{L}_\omega)\) from Subsection~\ref{sec:ch222}, each projector
block is an invariant subspace of the scalar Helmholtz symbol in the plane-wave
basis,
\medskip
Equations~\eqref{eq:ch2.spectral-projectors}--\eqref{eq:ch2.spectral-projector-algebra} are
\emph{formal} on the spectral completion \(\widehat{\mathcal{F}}_\Omega\): idempotency and
orthogonality are justified in the rigged Hilbert space treatment of Chapter~3 (\emph{Angular Momentum Operators and SO(3)})
Subsections~3.1.3--3.5.2 \cite{stratton1941}. Until that completion,
projector algebra is used only on admissible finite superpositions and on
\(\mathcal{R}_\omega[\Jimp]\), where outgoing selection is already enforced
\cite{harrington2001}. \begin{equation}
\label{eq:ch2.helmholtz-subspace-commutation}
\mathcal{L}_\omega\big(\Pi_{\mathrm{prop}}\mathbf{u}\big)
=
\Pi_{\mathrm{prop}}\big(\mathcal{L}_\omega\mathbf{u}\big),
\end{equation}
and similarly for \(\Pi_{\mathrm{ev}}\) and \(\Pi_{\mathrm{non}}\) when boundary
closure does not mix sheets \cite{harrington2001}. Equation
\eqref{eq:ch2.helmholtz-subspace-commutation} is the operator reason harmonic
solvers can treat propagating and evanescent content separately before radiation
selection; Sommerfeld outgoing radiation condition in Section~\ref{sec:ch24} acts only on
\(\Pi_{\mathrm{prop}}\) \cite{stratton1941}.
The decomposition
\begin{equation}
\label{eq:ch2.propagating-evanescent-decomposition}
\mathcal{S}_\omega[\Jimp]
=
\Pi_{\mathrm{prop}}\mathcal{S}_\omega[\Jimp]
+\Pi_{\mathrm{ev}}\mathcal{S}_\omega[\Jimp]
+\Pi_{\mathrm{non}}\mathcal{S}_\omega[\Jimp]
\end{equation}
holds when the impressed spectrum excites all three classes. Only
\(\Pi_{\mathrm{prop}}\mathcal{S}_\omega[\Jimp]\) is eligible for far-zone radiation
content after outgoing radiation condition; evanescent and non-propagating components remain in
reactive storage exchange \cite{stratton1941}. Energy norms from
Eq.~\eqref{eq:ch2.field-energy-norm} decompose orthogonally in the spectral basis
when the impressed angular spectrum has sufficient decay in \((k_x,k_y)\)
\cite{stratton1941}.
\begin{equation}
\label{eq:ch2.rad-projection-composition}
\Pi_{\mathrm{rad}}
=
\mathcal{O}_{\mathrm{out}}\circ\Pi_{\mathrm{prop}},
\end{equation}
where \(\mathcal{O}_{\mathrm{out}}\) enforces outgoing selection on propagating
content compatible with Eq.~\eqref{eq:ch1.radiation-trace-condition}
\cite{stratton1941}. Equation~\eqref{eq:ch2.rad-projection-composition} refines
\(\Pi_{\mathrm{rad}}\): radiative subspace \(\mathcal{F}_{\mathrm{rad}}\) is a
subset of \(\Pi_{\mathrm{prop}}\mathcal{F}_\Omega\), not of all of
\(\mathcal{F}_\Omega\) \cite{harrington2001}. Evanescent and non-propagating
components lie in the null space of far-zone radiation:
\begin{equation}
\label{eq:ch2.rad-null-on-evanescent}
\Pi_{\mathrm{rad}}\Pi_{\mathrm{ev}}=0,\qquad
\Pi_{\mathrm{rad}}\Pi_{\mathrm{non}}=0,
\end{equation}
while \(\Pi_{\mathrm{rad}}\Pi_{\mathrm{prop}}=\Pi_{\mathrm{rad}}\) after outgoing
closure \cite{stratton1941}. Equations~\eqref{eq:ch2.rad-null-on-evanescent} and
\eqref{eq:ch2.rad-projection-composition} restate, in projector form, the regime
logic of Section~\ref{sec:ch01.2.4}: reactive near fields may be large even when
\(\Pi_{\mathrm{rad}}\mathcal{S}_\omega[\Jimp]=\mathbf{0}\) \cite{harrington2001}.
Causality from Subsection~\ref{sec:ch223} ensures retarded support but does not,
by itself, promote evanescent content into \(\mathcal{F}_{\mathrm{rad}}\).
\medskip
For all \(r>r_{\min}(a,\lambda_0)\) and every transport, flux, near-to-far, and angular-spectrum
statement in Sections~\ref{sec:ch25}--\ref{sec:ch210},
\begin{equation*}
\Pi_{\mathrm{rad}}=\mathcal{O}_{\mathrm{out}}\circ\Pi_{\mathrm{prop}}
\quad\text{as in Eq.~\eqref{eq:ch2.rad-projection-composition}.}
\end{equation*}
Earlier uses of \(\Pi_{\mathrm{rad}}\) in Section~\ref{sec:ch20} are to be read through this
refinement \cite{stratton1941}. Transport qualification from Section~\ref{sec:ch01.2.4} enters as an operator criterion:
\begin{equation}
\label{eq:ch2.regime-operator-screen}
\mathcal{D}_{J,\mathrm{reg}}=1
\ \Longrightarrow\
\text{transport reports use}\ 
\Pi_{\mathrm{prop}}\mathcal{R}_\omega[\Jimp],
\ \text{not}\ \Pi_{\mathrm{ev}}\mathcal{S}_\omega[\Jimp],
\end{equation}
with \(\mathcal{D}_{J,\mathrm{reg}}\) from
Eq.~\eqref{eq:ch1.regime-qualified-decision} \cite{stratton1941}. Equation
\eqref{eq:ch2.regime-operator-screen} is the operator restatement of
Eq.~\eqref{eq:ch2.transport-qualified-radiation}; regime indices \(\chi(r)\) and
\(\chi_\Omega\) diagnose when Eq.~\eqref{eq:ch2.regime-operator-screen} fails on a
measurement shell \cite{stratton1941}.
Waveguide and cutoff structures illustrate non-propagating subspaces. For a uniform
guide with lateral eigenvalue \(\kappa_n\), mode \(n\) propagates only when
\begin{equation}
\label{eq:ch2.cutoff-condition}
k_0^2>\kappa_n^2
\quad\Longleftrightarrow\quad
\omega>\omega_{c,n}=\frac{c_0\kappa_n}{\sqrt{\varepsilon_r\mu_r}},
\end{equation}
with \(\omega_{c,n}\) the cutoff angular frequency \cite{stratton1941}. Below
cutoff, the mode lies in \(\mathcal{F}_{\mathrm{ev}}\) or
\(\mathcal{F}_{\mathrm{non}}\) depending on boundary termination; applying far-zone
transport formulas to such fields violates Eq.~\eqref{eq:ch2.regime-operator-screen}
as in the TE\(_{10}\) illustration of Section~\ref{sec:ch01.1.2}
\cite{harrington2001}. In a rectangular guide of width \(a\), lateral eigenvalues
\(\kappa_{m,n}=\sqrt{(m\pi/a)^2+(n\pi/b)^2}\) replace \(k_t\) in
Eq.~\eqref{eq:ch2.cutoff-condition}; each mode index \((m,n)\) switches from
\(\mathcal{F}_{\mathrm{non}}\) to \(\mathcal{F}_{\mathrm{prop}}\) only when
\(\omega\) exceeds its cutoff \cite{stratton1941}. Near cutoff the propagating
group velocity vanishes and stored reactive energy dominates even though the mode
has crossed into \(\Pi_{\mathrm{prop}}\); Eq.~\eqref{eq:ch1.radiative-regime-condition}
remains the far-zone transport criterion, not Eq.~\eqref{eq:ch2.cutoff-condition} alone
\cite{harrington2001}.
Angular-spectrum expansion: expand \(\E(\mathbf{r})\) on \(z=\mathrm{const}\) as
\begin{equation}
\label{eq:ch2.angular-spectrum}
\E(\mathbf{r})
=
\int\!\!\int \widetilde{\E}(k_x,k_y)\,
e^{-\jj(k_x x+k_y y+k_z(k_t) z)}\,dk_x\,dk_y,
\end{equation}
with \(k_z(k_t)\) from Eq.~\eqref{eq:ch2.longitudinal-wavenumber-split}
\cite{stratton1941}. Equation~\eqref{eq:ch2.angular-spectrum} is the
Chapter~\ref{ch:02} spectral template for Section~\ref{sec:ch29}; full radiation
operator action on angular spectra is developed there. Subsection~\ref{sec:ch224} only
defines the propagating/evanescent partition induced by
Eq.~\eqref{eq:ch2.dispersion-relation} \cite{harrington2001}.
\begin{figure}[t]
\centering
\ChOneFigBlock{%
\begin{tikzpicture}[ch1fig, node distance=5mm and 7mm]
 \node[ch1box] (s) {$\mathcal{S}_\omega[\Jimp]$};
 \node[ch1box, right=7mm of s] (pp) {$\Pi_{\mathrm{prop}}$};
 \node[ch1box, right=7mm of pp] (pr) {$\Pi_{\mathrm{rad}}$}; \node[ch1box, below=5mm of s] (pe) {$\Pi_{\mathrm{ev}}$};
 \node[ch1box, below=5mm of pe] (pn) {$\Pi_{\mathrm{non}}$};
 \draw[ch1flow] (s) -- (pp) -- (pr);
 \draw[ch1flow] (s.south) -- (s.south |- pe.north);
 \draw[ch1flow] (s.south |- pe.north) -- (pe.north);
 \draw[ch1flow] (s.east) -- (s.east |- pn.west);
 \draw[ch1flow] (s.east |- pn.west) -- (pn.west);
\end{tikzpicture}%
}
\caption{Spectral criterion: only \(\Pi_{\mathrm{prop}}\) content passes to
\(\Pi_{\mathrm{rad}}\); evanescent and non-propagating components remain reactive.}
\label{fig:ch2.spectral-radiation-criterion}
\end{figure}
Figure~\ref{fig:ch2.source-radiation-sequence}. Causality from
Subsection~\ref{sec:ch223} does not automatically place evanescent energy in
\(\mathcal{F}_{\mathrm{rad}}\); spectral and regime qualifications are additional
\cite{stratton1941}. Near-zone reactive fields decompose into large evanescent and non-propagating
fractions even when \(\mathcal{S}_\omega[\Jimp]\) is exact. As \(r\to\infty\),
\(\Pi_{\mathrm{prop}}\) content dominates \(P_{\mathrm{rad}}\) when
Eq.~\eqref{eq:ch1.radiative-regime-condition} holds \cite{stratton1941}. Operator
norms on subspaces satisfy
\begin{equation}
\label{eq:ch2.subspace-energy-inequality}
\|\Pi_{\mathrm{ev}}\mathcal{S}_\omega[\Jimp]\|_{\mathcal{U}}
\gg
\|\Pi_{\mathrm{prop}}\mathcal{S}_\omega[\Jimp]\|_{\mathcal{U}}
\quad\text{on near shells,}
\end{equation}
while the reverse inequality governs far-zone radiative observables
\cite{harrington2001}. Subsection~\ref{sec:ch224} does not compute
shell integrals; it supplies the spectral vocabulary for
Section~\ref{sec:ch252} and flux operators in Section~\ref{sec:ch25}.
With Subsections~\ref{sec:ch221}--\ref{sec:ch224} complete, Section~\ref{sec:ch22}
has connected time-domain evolution, harmonic Maxwell and Helmholtz operators,
causal equivalence, and propagating--evanescent spectral split on the function spaces
of Section~\ref{sec:ch21}. Green-function inversion in Section~\ref{sec:ch23}
realizes \(\mathcal{L}_\omega^{-1}\) and \(\mathcal{G}_{E,\omega}\) with retarded
causality; Sommerfeld selection in Section~\ref{sec:ch24} fixes outgoing branches for
\(\Pi_{\mathrm{prop}}\) and \(\Pi_{\mathrm{rad}}\); radiation integrals in
Section~\ref{sec:ch25} meter flux on the propagating subspace qualified by
Eq.~\eqref{eq:ch2.regime-operator-screen} \cite{stratton1941}.
